%%%%%%%%%%%%%%%%%%%%%%%%%%%%%%%%%%%%%%%%%%%%%%%%%%%%%%%%%%%%%%%%%%%%%%%%%%%%%%%%

\pagestyle{empty}
\begin{abstract}
  La digitalización de los clubes deportivos es una necesidad creciente en el
  contexto del auge del pádel en España. Sin embargo, la mayoría de los centros
  deportivos siguen gestionando sus reservas de forma manual o con herramientas
  genéricas que no cubren sus necesidades específicas. PadelCenter es una
  aplicación web diseñada para resolver este problema, ofreciendo una plataforma
  integral para la gestión de reservas de pistas de pádel y la organización de
  torneos.

  El sistema se ha construido siguiendo principios de arquitectura hexagonal,
  separando el dominio de negocio de los detalles tecnológicos, lo que facilita
  el mantenimiento y la evolución del sistema. El backend está desarrollado en
  Java~21 con Spring Boot~3, empleando un enfoque \textit{contract-first} basado
  en OpenAPI~3 para definir el contrato de la API antes de implementar cualquier
  lógica. El frontend es una aplicación Next.js~14 con App Router, TypeScript y
  React Query, con clientes HTTP generados automáticamente a partir de la
  especificación OpenAPI mediante Orval.

  Las funcionalidades principales incluyen autenticación segura mediante JWT, la
  gestión de centros deportivos y pistas, la creación y consulta de reservas con
  control de conflictos, y la organización de torneos con generación automática
  de emparejamientos. La estrategia de pruebas abarca tests unitarios con
  Mockito, tests de integración con Testcontainers y tests de arquitectura con
  ArchUnit.

  El resultado es un sistema funcional, bien estructurado y extensible, que
  demuestra la aplicabilidad de las prácticas modernas de ingeniería del software
  en el ámbito de la gestión deportiva.
\end{abstract}

\clearpage
\pagestyle{empty}
\begin{description}
\item [\palabraschaveprincipal:] \mbox{} \\[-20pt]
  \begin{itemize}
    \item Arquitectura hexagonal
    \item Spring Boot
    \item Next.js
    \item API REST
    \item OpenAPI
    \item Gestión deportiva
    \item Reservas
    \item Torneos
    \item Java
    \item TypeScript
  \end{itemize}
\end{description}

\pagestyle{fancy}

%%%%%%%%%%%%%%%%%%%%%%%%%%%%%%%%%%%%%%%%%%%%%%%%%%%%%%%%%%%%%%%%%%%%%%%%%%%%%%%%
